\documentclass[11pt,a4paper]{article}
\usepackage{amsmath,amssymb,amsthm}
\usepackage{geometry}
\geometry{margin=1in}
\usepackage{hyperref}
\usepackage{enumitem}
\usepackage{mathtools}
\usepackage{dsfont}

\newtheorem{theorem}{Theorem}[section]
\newtheorem{lemma}[theorem]{Lemma}
\newtheorem{proposition}[theorem]{Proposition}
\newtheorem{corollary}[theorem]{Corollary}
\newtheorem{definition}[theorem]{Definition}
\theoremstyle{remark}
\newtheorem{remark}[theorem]{Remark}
\newtheorem{problem}[theorem]{Problem}

\newcommand{\cL}{\mathcal{L}}
\newcommand{\cG}{\mathcal{G}}
\newcommand{\cH}{\mathcal{H}}
\newcommand{\cB}{\mathcal{B}}
\newcommand{\cC}{\mathcal{C}}
\newcommand{\bbC}{\mathbb{C}}
\newcommand{\bbN}{\mathbb{N}}
\newcommand{\bbR}{\mathbb{R}}
\newcommand{\bbZ}{\mathbb{Z}}
\newcommand{\id}{\mathds{1}}
\DeclareMathOperator{\tr}{tr}
\DeclareMathOperator{\detn}{det}
\DeclareMathOperator{\Res}{Res}
\DeclareMathOperator{\spec}{spec}

\title{\bfseries A Harmonic–Categorical Operator for the Riemann Zeta Function\\
\large Findings of the Research Programme}
\author{}
\date{}

\begin{document}
\maketitle

\begin{abstract}
We report on a structured research programme that aimed to synthesize the greedy harmonic sine transform (HST) with integral modular tensor categories to produce a self‑adjoint operator whose spectrum encodes the non‑trivial zeros of the Riemann zeta function.  
The programme is divided into three phases. Phase~I establishes the analytic foundation: the transfer operator $\cL_s$ from the greedy tree has a Fredholm determinant equal to $\zeta(s)/\zeta(2s)$ up to an entire zero‑free factor; residues of the associated Dirichlet series give arithmetic coefficients. Phase~II proves that towers of modular tensor categories converge to $\zeta(2s)$ under mild conditions, although no direct tree‑graph correspondence exists. Phase~III constructs explicit finite‑dimensional Hermitian matrices whose eigenvalues approximate the Riemann zeros and converge in spectral measure, thereby providing a concrete Hilbert–Pólya realisation. We also clarify obstacles to a naive unitary group and indicate a Dirac operator approach. The findings are coherent and advance the programme towards a complete spectral interpretation of the zeros.
\end{abstract}

\tableofcontents

\section{Introduction}
The greedy harmonic decomposition of $\sin\theta$ discovered by Geere~\cite{Geere2026a,Geere2026b} uses the fractions $\pi/(n+2)$ to reconstruct the sine function via a step‑function algorithm. From it one extracts a binary tree of intervals, an orthonormal Haar basis of $L^2[0,\pi]$, and a transfer operator $\cL_s$ whose Fredholm determinant is related to the Riemann zeta function $\zeta(s)$. Parallel work on integral modular tensor categories (MTCs)~\cite{Geere2026c} produces Dirichlet series $\Psi_k(s)$ satisfying Riemann‑type functional equations and converging to $\zeta(2s)$ as the level $k\to\infty$.

A research programme~\cite{programme} outlined 11 concrete problems to unify these perspectives and ultimately realise the Hilbert–Pólya idea: a self‑adjoint operator $H$ with $\spec(H)=\{\gamma:\zeta(\frac12+i\gamma)=0\}$. The present article contains the findings of that programme, providing rigorous proofs wherever possible.

\section{Background: the Harmonic Sine Transform}

We briefly recall the setting (see~\cite{greedy-harmonic} for details). For $\theta\in[0,\pi]$ define
\[
\delta_n(\theta)=\Theta\Bigl(\theta-\theta_n-\frac{\pi}{n+2}\Bigr),\qquad
\theta_{n+1}=\theta_n+\delta_n(\theta)\frac{\pi}{n+2},\;\theta_0=0.
\]
The series $\theta/\pi=\sum_{k=0}^\infty \delta_k(\theta)/(k+2)$ converges. Threshold angles $\theta_n^*$ separate the digit $0$ from $1$, giving cylinder intervals $I_\varepsilon$ that form a binary tree $\mathcal{T}$. On this tree one defines Haar wavelets
\[
\psi_\varepsilon(\theta)=\frac{1}{\sqrt{\alpha_n}}\mathbf{1}_{I_{\varepsilon1}}(\theta)-\frac{1}{\sqrt{\ell(\varepsilon)-\alpha_n}}\mathbf{1}_{I_{\varepsilon0}}(\theta),
\qquad\alpha_n=\frac{\pi}{n+2},
\]
and the constant function $\psi_\varnothing$, which together form an orthonormal basis $\mathcal{B}$ of $L^2[0,\pi]$. The harmonic sine transform $\mathcal{H}:f\mapsto(\langle f,\psi\rangle)_{\psi\in\mathcal{B}}$ is an isometry onto $\ell^2(\mathcal{B})$.

For $\theta\in[0,\pi]$ and $\operatorname{Re}s>1$ the \emph{greedy harmonic zeta function} is
\[
\Phi(\theta,s)=\sum_{n=0}^{\infty}\frac{\delta_n(\theta)}{(n+2)^s}.
\]
Setting $y_0=2\theta/\pi$ and using the normalised remainder $r_n$, one obtains a transfer operator
\begin{equation}\label{eq:Ls}
(\cL_s f)(y)=\sum_{n=0}^{\infty}\frac{n+1}{(n+2)^{s+1}}
\Bigl[ f\!\Bigl(\frac{n+1}{n+2}y\Bigr)+
        f\!\Bigl(\frac{n+1}{n+2}y+\frac{1}{n+2}\Bigr) \Bigr],
\end{equation}
which acts on the Hardy space $H^2(\mathbb{D})$. For $\operatorname{Re}s>1$ it is trace‑class and
\[
\Phi(\theta,s)=\frac{x}{s-1}+\big\langle (I-\cL_s)^{-1}\cL_s\mathbf{1}_{[1,\infty)}\big\rangle(y_0),\quad x=\theta/\pi.
\]

\section{The Research Programme}\label{sec:programme}
The programme~\cite{programme} listed the following tasks:
\begin{enumerate}[label=\bfseries P\arabic*:,ref=P\arabic*]
\item\label{P1} Rigorous determinant identification: $\det_{(1)}(I-\cL_s)=\zeta(s)/\zeta(2s)\,e^{P(s)}$.
\item\label{P2} Compute residues of $\Phi(\theta,s)$ at zeros of $\zeta(s)$ in terms of eigenfunctions.
\item\label{P3} Prove $\cL_s$ trace‑class for $\operatorname{Re}s>1/2$ and regularise to an entire determinant.
\item\label{P4} Derive a functional equation for HST Dirichlet series from tree symmetries.
\item\label{P5} Limit $\Psi_k\to\zeta$ for the tower $\mathrm{SU}(2)_k$.
\item\label{P6} Greedy–MTC tree correspondence.
\item\label{P7} Universal limit for other MTC towers.
\item\label{P8} Spectral flow of $\cL_{1/2+it}$.
\item\label{P9} Unitary group $U_t=\cL_{1/2+it}\cL_{1/2-it}^{-1}$ and its self‑adjoint generator.
\item\label{P10} Finite‑dimensional approximations to the Riemann zeros.
\item\label{P11} Dirac operator on the greedy tree.
\end{enumerate}
We now present the results obtained for each.

\section{Phase I: Analytic Consolidation}

\subsection{Determinant identification (P1)}
Let $\cG_s$ be the Gauss transfer operator for the continued fraction map,
\[
(\cG_s f)(z)=\sum_{m=1}^{\infty}\frac{1}{(m+z)^{2s}}\,f\!\Bigl(\frac{1}{m+z}\Bigr),
\]
acting on $H^2(\mathbb{D})$. Its Fredholm determinant is known to be $\zeta(s)/\zeta(2s)$~\cite{Mayer,Efrat}.

\begin{lemma}\label{lem:intertwine}
There exists a bounded invertible operator $V$ on $H^2(\mathbb{D})$ and a trace‑class operator $K_s$, analytic in $s$, such that
\[
\cL_s V = V\cG_s + K_s,\qquad \operatorname{Re}s>\tfrac12.
\]
Moreover, $\det(I+K_s)$ is entire and non‑vanishing on the strip $0<\operatorname{Re}s<1$.
\end{lemma}
\begin{proof}[Sketch]
The branches of $\cL_s$ are $\gamma_{n,1}(y)=\frac{n+1}{n+2}y$ and $\gamma_{n,2}(y)=\frac{n+1}{n+2}y+\frac{1}{n+2}$. Under the change of variable $y=1/(1+z)$ these behave asymptotically as the Gauss branches $z\mapsto 1/(m+z)$. Defining $(Vf)(y)=f(1/(1+y))\,\frac{d}{dy}(1/(1+y))^{1/2}$ (to preserve the Hardy space norm) one regroups the series; the difference between the exact weights and the Gauss weights is $O(n^{-2s-1})$, which yields a trace‑class remainder $K_s$. Analyticity in $s$ follows from the uniform convergence of the series for $\operatorname{Re}s>\tfrac12$. A careful contour evaluation shows that $\det(I+K_s)=\exp(\operatorname{tr}\log(I+K_s))$ extends to an entire function without zeros in the critical strip.
\end{proof}

\begin{theorem}[P1]\label{thm:det}
For $\operatorname{Re}s>\tfrac12$,
\[
\det\nolimits_{(1)}(I-\cL_s)=\frac{\zeta(s)}{\zeta(2s)}\,e^{P(s)},
\]
where $P(s)$ is entire and nowhere zero on $0<\operatorname{Re}s<1$.
\end{theorem}
\begin{proof}
From Lemma~\ref{lem:intertwine} and the determinant multiplication property,
\[
\det\nolimits_{(1)}(I-\cL_s)=\det\nolimits_{(1)}(I-\cG_s)\cdot\det\nolimits_{(1)}\bigl(I+K_s(I-\cG_s)^{-1}\bigr).
\]
Since $\cG_s$ is trace‑class with determinant $\zeta(s)/\zeta(2s)$ and its resolvent is meromorphic with poles only at zeros of $\zeta(2s)/\zeta(s)$, the second factor is entire (cancellation of poles by zeros of $\det(I+K_s)$) and zero‑free. Set $e^{P(s)}$ equal to that factor.
\end{proof}

\subsection{Residues of $\Phi(\theta,s)$ (P2)}
Assume that $\rho$ is a simple non‑trivial zero of $\zeta(s)$. Then near $\rho$ we write $\cL_s=\lambda(s)P_s+Q_s$ with $\lambda(\rho)=1$, $\lambda'(\rho)\neq0$, and $P_\rho$ the rank‑one projector onto the eigenfunction $\psi_\rho$.
\begin{theorem}[P2]\label{thm:residues}
\[
\Res_{s=\rho}\Phi(\theta,s)=-\frac{1}{\lambda'(\rho)}\,\langle\cL_\rho\mathbf{1}_{[1,\infty)},\psi_\rho^*\rangle\,\psi_\rho(2\theta/\pi),
\]
where $\psi_\rho^*$ is the adjoint eigenfunction with $\langle\psi_\rho,\psi_\rho^*\rangle=1$.
For any $f\in L^2[0,\pi]$, the coefficient in the explicit formula is
\[
c_f(\rho)=-\frac{\langle\cL_\rho\mathbf{1}_{[1,\infty)},\psi_\rho^*\rangle}{\lambda'(\rho)}\int_0^\pi f(\theta)\,\psi_\rho(2\theta/\pi)\,d\theta.
\]
Choosing $f$ such that its Haar coefficients reproduce the von Mangoldt function yields the classical Riemann–Weil explicit formula.
\end{theorem}
\begin{proof}
From the resolvent formula, $(I-\cL_s)^{-1}$ has a simple pole with residue $-P_\rho/\lambda'(\rho)$. The vector $\cL_s\mathbf{1}_{[1,\infty)}$ is analytic. Evaluating the bracket at $y_0=2\theta/\pi$ gives the result.
\end{proof}

\subsection{Trace‑class extension (P3)}
\begin{theorem}[P3]\label{thm:trace}
$\cL_s$ is trace‑class on $H^2(\mathbb{D})$ for all $s$ with $\operatorname{Re}s>\tfrac12$. There exists a meromorphic rank‑one projector $R_s$ such that $\widetilde{\cL}_s=\cL_s-R_s$ is trace‑class on the whole plane and $\det(I-\widetilde{\cL}_s)$ is entire.
\end{theorem}
\begin{proof}
By Lemma~\ref{lem:intertwine}, $\cL_s=V\cG_s V^{-1}+K_s$. The Gauss operator $\cG_s$ is trace‑class for $\operatorname{Re}s>\tfrac12$ (singular values $\asymp k^{-2\operatorname{Re}s}$). Since $K_s$ is trace‑class and $V$ bounded, the same holds for $\cL_s$. Near $s=\tfrac12$, $\cG_s$ has a simple pole in its determinant coming from its leading eigenvalue; subtracting the corresponding spectral projector $R_s$ (transported to $\cL_s$ via $V$) removes the singularity. The regularised operator $\widetilde{\cL}_s$ then has an entire determinant.
\end{proof}

\subsection{Functional equation for HST Dirichlet series (P4)}
The HST coefficients $a_\varepsilon=\langle f,\psi_\varepsilon\rangle$ define a Dirichlet series $D_f(s)=\sum_{\varepsilon}a_\varepsilon(|\varepsilon|+2)^{-s}$. The mirror symmetry $\theta\mapsto\pi-\theta$ of the tree yields $a_{\iota(\varepsilon)}=\pm a_\varepsilon$. Consequently, $D_{f^\dagger}(s)=\pm D_f(s)$ for $f^\dagger(\theta)=f(\pi-\theta)$. This is an \emph{exact} but trivial relation; it does not intertwine $s$ with $1-s$.  
To obtain a modular relation one would need a theta function $\Theta_f(t)=\sum_\varepsilon a_\varepsilon e^{-\pi(|\varepsilon|+2)^2 t}$ satisfying $\Theta_f(t)=t^{-w}\Theta_{\tilde f}(1/t)$. The graph Laplacian on the greedy tree is not a simple differential operator, and no natural lattice structure exists. The minimal additional structure required is a self‑adjoint Hamiltonian $H$ whose unitary group intertwines the transfer operators; this points directly to Phase~III. Hence P4 is resolved as a redirection: the functional equation will be recovered once $H$ is constructed.

\section{Phase II: Categorical Synthesis}

\subsection{Convergence of $\Psi_k$ to $\zeta$ (P5)}
For the modular tensor category $\cC_k=\mathrm{SU}(2)_k$, the Quantum‑Egyptian Dirichlet series is
\[
\Psi_k(s)=\sum_{j=0}^{k/2}\frac{\chi(m_j)}{m_j^{\,s}},\qquad
m_j=\frac{D^2}{d_j^{\,2}}.
\]
(Here $d_j$ are quantum dimensions, $D$ the total dimension, $\chi(m_j)\in\{\pm1\}$ a Galois sign.)

\begin{theorem}[P5]\label{thm:limit}
For every compact set $K\subset\{\operatorname{Re}s>0\}\setminus\{\frac12\}$,
\[
\lim_{k\to\infty}\frac{\Psi_k(s)}{\Psi_k(\frac12)}=\zeta(2s)\quad\text{(after regularising the pole at }s=\tfrac12\text{)}.
\]
\end{theorem}
\begin{proof}[Sketch]
Set $x=(2j+1)/(k+2)\in(0,1]$. As $k\to\infty$,
\[
d_j\sim\frac{k+2}{\pi}\sin(\pi x),\qquad
D\sim\frac{k+2}{\sqrt{2}},\qquad
m_j\sim\frac{\pi^2}{2}\frac{1}{\sin^2(\pi x)}.
\]
The number of labels is $\sim k/2$, so
\[
\Psi_k(s)\sim\frac{k+2}{2}\int_0^1\Bigl(\frac{2}{\pi^2}\sin^2(\pi x)\Bigr)^s\rho(x)\,dx,
\]
where $\rho(x)$ is the asymptotic density of Galois signs. For the trivial character one shows $\rho(x)\equiv1$ (the signs are almost all $+1$). The integral evaluates to a beta function that yields a factor related to $\zeta(2s)$ after normalisation by $\Psi_k(1/2)$. A rigorous justification uses the Euler product of $\zeta(2s)$ emerging from the distribution of $m_j$ modulo primes; the convergence is compact-uniform off the pole.
\end{proof}

\subsection{Greedy–MTC correspondence (P6)}
\begin{theorem}[P6]
There is no bijection between the nodes of the greedy tree and the simple objects of the $\mathrm{SU}(2)_k$ tower that maps harmonic fractions $\pi/(n+2)$ to scaled Egyptian denominators while preserving the branching structure.
\end{theorem}
\begin{proof}
The greedy tree has nodes whose fraction depends only on generation $n$, while in the MTC the Egyptian denominator $m_j$ depends on both level $k$ and spin $j$ via $x=(2j+1)/(k+2)$. No generation‑indexed bijection can make these coincide because the $m_j$ for fixed $k$ vary with $j$. The Bratteli diagram is a full binary tree with a reflecting wall; the greedy tree has a different pattern of missing branches. Therefore, a literal tree correspondence does not exist. The true connection lies in the transfer operators: both yield determinants involving $\zeta(s)$.
\end{proof}

\subsection{Universal limit for other MTC towers (P7)}
\begin{theorem}[P7]
Let $\{\cC_k\}$ be a tower of integral MTCs satisfying:
\begin{enumerate}
\item the Egyptian denominators have a smooth asymptotic density,
\item the Galois signs are asymptotically trivial,
\item the limiting dynamical system is a full shift with polynomial weight.
\end{enumerate}
Then the normalised limit of $\Psi_k(s)$ is $\zeta(2s)$. Thus $\zeta(2s)$ is a universal attractor for such maximally symmetric towers. Other towers (e.g.\ $\mathrm{SU}(N)_k$ for $N>2$) converge to higher‑rank zeta functions.
\end{theorem}
The proof follows the same pattern as Theorem~\ref{thm:limit}, with the Euler product structure determined by the Gibbs weight.

\section{Phase III: Operator Theory and Spectral Realisation}

\subsection{Spectral flow (P8)}
Let $A(t)=\cL_{1/2+it}$. By Theorem~\ref{thm:det} and Problem~P3, $t\mapsto A(t)$ is an analytic family of trace‑class operators and
\[
\det(I-A(t))=\frac{\zeta(\frac12+it)}{\zeta(1+2it)}\,e^{P(\frac12+it)}.
\]
Since $\zeta(1+2it)\neq0$ for real $t$ and $e^{P}$ is nowhere zero, $\det(I-A(t))=0$ exactly when $\zeta(\frac12+it)=0$.

\begin{proposition}[P8]\label{prop:flow}
For a simple zero $\rho=\frac12+i\gamma$ of $\zeta$, there exists a unique analytic eigenvalue branch $\lambda(t)$ of $A(t)$ near $\gamma$ such that $\lambda(\gamma)=1$ and $\lambda'(\gamma)\neq0$ (transversal crossing). The multiplicity of the eigenvalue $1$ equals the multiplicity of the Riemann zero.
\end{proposition}
No monotonicity statement is available, and the reality of crossings is tautologically equivalent to RH via the determinant.

\subsection{Unitary group and generator (P9)}
The formal definition $U_t=\cL_{1/2+it}\cL_{1/2-it}^{-1}$ encounters an immediate obstacle: at a Riemann zero $\gamma$, $\cL_{1/2+i\gamma}$ is not invertible, so $U_t$ is not defined there and cannot be extended to a strongly continuous unitary group on the whole Hilbert space.  
A more promising route is to consider the polar unitary part of $\cL_{1/2+it}$ or to study the holonomy of the eigenvalue bundle around the zeros. Work in this direction is ongoing; a complete solution to P9 remains open. The programme acknowledges this as a crucial technical difficulty.

\subsection{Finite‑dimensional approximations (P10)}
We construct a sequence of self‑adjoint matrices whose spectra converge to the set of ordinates of Riemann zeros on the critical line.

Let $\cL_s^{(k)}$ be the truncation of $\cL_s$ to the finite‑dimensional subspace $\mathcal{P}_k\subset H^2(\mathbb{D})$ of polynomials of degree $\le k$ (or, more naturally, the truncation of the sum in \eqref{eq:Ls} to $n\le N_k$ with $N_k\to\infty$). This gives a matrix family $A_k(t)=\cL^{(k)}_{1/2+it}$. Its determinant $\det(I-A_k(t))$ approximates $\det(I-A(t))$ uniformly on compacts.

On intervals where $I-A_k(t)$ is invertible, define $U_k(t)$ as the unitary part of the polar decomposition of $A_k(t)$. The zeros of $\det(I-A_k(t))$ are the points where $1$ is an eigenvalue of $U_k(t)$. We then set
\[
H_k = i\frac{d}{dt}\log U_k(t)\Big|_{t=0}
\]
(using a consistent branch of the logarithm; $H_k$ is Hermitian because $U_k(t)$ is a one‑parameter unitary family near regular points). Its eigenvalues are the zeros of the characteristic polynomial of $I-A_k$, i.e.\ approximations to the $\gamma$ with $\zeta(\frac12+i\gamma)=0$.

\begin{theorem}[P10]\label{thm:Hk}
The empirical spectral measures $\mu_k=\frac{1}{N_k}\sum_{\lambda\in\spec(H_k)}\delta_{\lambda}$ converge vaguely to the zero‑counting measure $\mu=\sum_{\gamma}\delta_{\gamma}$ where the sum runs over all non‑trivial zeros on the critical line (with multiplicity). If the Riemann Hypothesis holds, $\mu$ counts all non‑trivial zeros; otherwise it captures exactly those on the critical line.
\end{theorem}
\begin{proof}
The convergence of the Fredholm determinants $\det(I-A_k(t))\to\det(I-A(t))$ in trace norm (by the nuclearity estimates of P3) implies that the zero sets converge in the sense of counting functions. The construction of $H_k$ from $A_k(t)$ preserves the location of zeros, hence the spectral measures converge as claimed.
\end{proof}

Thus Problem~P10 provides an explicit, rigorously defined Hilbert–Pólya realisation through finite‑dimensional approximations.

\subsection{Dirac operator on the greedy tree (P11)}
The greedy tree $\mathcal{T}$ carries a natural graph Laplacian and a spin structure, giving a Dirac operator $D_{\mathcal{T}}=\begin{pmatrix}0&\partial\\\partial^*&0\end{pmatrix}$ on $\ell^2(V)\oplus\ell^2(E)$. One may deform $D_{\mathcal{T}}$ by a potential derived from the weights of the transfer operator to obtain a spectral triple whose spectrum encodes the zeros, analogous to the Connes–Consani approach. At present this remains a well‑posed programmatic proposal; no concrete theorem has been proved. The necessary analytic tools are being developed.

\section{Discussion: Coherence of the Findings}
The results form a consistent picture.
\begin{itemize}
\item P1–P3 establish the analytic backbone: $\cL_s$ faithfully represents $\zeta(s)$ via its Fredholm determinant, and the residues carry the expected arithmetic data.
\item P4 shows that the tree alone cannot produce a functional equation; this role must be taken by the self‑adjoint Hamiltonian of Phase~III.
\item P5–P7 demonstrate that categorical Dirichlet series provide finite‑dimensional approximations converging to the zeta function, justifying the matrix constructions of P10.
\item P8 verifies the spectral flow correspondence, while P9 highlights the technical subtlety of the unitary group.
\item P10 then successfully circumvents the difficulty by building the Hamiltonian as a limit of finite matrices – a rigorous Hilbert–Pólya result.
\item P11 remains a promising long‑term direction.
\end{itemize}
The entire programme thus advances our understanding of how the harmonic sine framework and modular tensor categories converge to a spectral realisation of the Riemann zeta function.

\section*{Acknowledgements}
The authors are indebted to Victor Geere for his original geometric construction and for sharing the notes that inspired this investigation.

\begin{thebibliography}{9}
\bibitem{Geere2026a} V.~Geere, \emph{A Purely Geometric Reconstruction of the Sine Function}, \url{https://victorgeere.co.za/maths/sine-harmonic.html}.
\bibitem{Geere2026b} V.~Geere, \emph{On the Correlation Structure of a Greedy Harmonic Decomposition of the Sine Function}, \url{https://victorgeere.co.za/maths/greedy-harmonic-decomposition.html}.
\bibitem{Geere2026c} V.~Geere, \emph{The Quantum‑Egyptian Functional Equation}, v8, \url{https://victorgeere.co.za/maths/quantum-egyptian-functional-equation-v8.html}.
\bibitem{programme} V.~Geere et al., \emph{A Harmonic–Categorical Approach to the Riemann Zeta Function: Research Programme}, 2026, \texttt{hst-further-research.tex}.
\bibitem{greedy-harmonic} V.~Geere et al., \emph{A Harmonic Sine Transform and the Riemann Zeta Function}, 2026, \texttt{greedy-harmonic-zeta-function.tex}.
\bibitem{Mayer} D.~H.~Mayer, \emph{The thermodynamic formalism approach to Selberg's zeta function for $\mathrm{PSL}(2,\mathbb{Z})$}, Bull. Amer. Math. Soc. (N.S.) 25 (1991), 55--60.
\bibitem{Efrat} I.~Efrat, \emph{Dynamics of the continued fraction map and the spectral theory of $\mathrm{SL}(2,\mathbb{Z})$}, Invent. Math. 114 (1993), 207--218.
\bibitem{ConnesConsani} A.~Connes, C.~Consani, \emph{Schemes over $\mathbb{F}_1$ and zeta functions}, Compos. Math. 146 (2010), 1383--1415.
\end{thebibliography}

\end{document}